% Logical positivism ask for philosophical care to be taken at just this point. Measuring the angles between light beams helo us to determine the angle sum property for physical triangles only if light travels in straight lines. How can you determine this? you could see if light can be beamed along the edge of a straight edge. But this again presuppose that the straightefge is straight, and how do you determine that Eventually, say the logical positivists, building on their understanding of the work of Einstein and Heinri Poincaré you get to a point where you cannot empirically check for physical strightness bu rather must take a conventional choice about what are you going to count as straight. Within the parlance employed by blumberg anf Feigl, which they took from the work of reichenbach, the scientist must take a conventional choice pf an axio;m or definition of coordination.


%When we encounter such “information destroying” processes, then,
%we should doubt our capacity to uncover the past. Dick Lewontin makes
%similar complaints about attempts to reconstruct the evolution of human
%cognition: “History, and evolution is a form of history, simply does not
%leave sufficient traces, especially when it is the forces that are at issue. Form
%and even behavior may leave fossil remains, but forces like Natural Selec-
%tion do not” (Lewontin 1998, 132).
%


If phrased in the past tense, these statements
seem trivial—all these things have indeed existed. If historical realism was only
about this, we would all be realists. For even a strong historical anti-realist about
“the past” will surely accept that Caesar, Cleopatra, and Sima Qian have existed.

However, the problem of historical realism may not be about the necessary
presuppositions of historical inquiry; instead, it may be about what must be the
case about the world for historical knowledge to be possible, independent of what
historians presuppose or believe. In this case, we have two different possibilities:
(1) if the problem of historical realism turns not on existence but on accessibil-
ity/manipulability, then it would remain separate from the presentism-eternalism
debate and lead us to epistemological issues, or (2) if the problem of historical
realism does turn on the existence of past entities or processes, then it is reducible
to the presentism-eternalism debate. Case (1) considers that interlocutors may re-
alize that what drives their disagreement is not the problem of existence per se;
rather, it concerns the use of “exist” and the lack of accessibility or manipulabil-
ity of past entities or processes. In this understanding, historical anti-realism may
need not lead to presentism. Rather, given that our human frames of reference are
subject to similar conditions under general relativity (we are all on planet Earth,
subject to similar physical forces, and possess similar cognitive capacities), his-
torical anti-realism turns into the claim that we are bound not to have any mean-
ingful way of accessing or interacting with past entities or processes, regardless
of their metaphysical-existential status. This would offer a way to differentiate the
problem of historical realism from the presentism-eternalism debate.

The argument isj pretty easy. If you consider history as a science and ai tell you must fo exxactly that, then our debates on realism anti-realism become trivializes cause no jsignnle individual would deny the existetnce of at least some of the founders of their  own discipline. If not, then you're stupid.


There is another x that can provide a frame to the problem of historical real-
ism: narrative. Of course, our main worry is not whether narrative as a discursive
form in general exists independently but whether our narratives about past en-
tities or processes somehow grasp a connective structure (both synchronic and
diachronic) that is an independent feature of reality. Then, our question model
can be expressed as follows: Do these connective structures exist independently?
Narrative realism is the claim that, yes, this complex configuration of relation-
ships between past entities or processes is something that exists independently—
for instance, this is one of the claims in Currie and Swaim’s minimal realism

JOÃO OHARA 15
from problems of evidence, method, and conceptual structure. But there are also
the problems of interpretation and language, insofar as the descriptions of enti-
ties and processes that interest historians are often not the mere description of
physical states or properties. And, finally, just as in our previous exploration of
metaphysics, here, too, the definition of the target can make an important differ-
ence.
Let us start from this last point. Take the examples of past “events” and “social
structures”: each of these may lead us to different reactions. 30 General skepti-
cism about (past) events as a kind is difficult to even imagine. Take an instance of
an event such as, “on 24 August 1954, Brazilian president Getúlio Vargas com-
mitted suicide at the presidential palace, the Catete Palace, in Rio de Janeiro.”
It left plenty of material remains, both direct (the gun, the body, the president’s
clothes) and indirect (a suicide letter, journalistic coverage, other kinds of docu-
mentary records), that could be—and, in fact, were—studied in research projects
concerned with many different questions. There may be other events for which it
is difficult or impossible to provide plausible evidential support, but that is not a
problem of the “event” kind; rather, it is a problem of the specific instance under
consideration. We are skeptical of the statement that “Big Foot broke my window
last night” not because we are skeptical about events but because we are skepti-
cal about this instance of an alleged event. Believing that it is possible to achieve
historical knowledge of events is perfectly compatible with the existence of either
unprovable events or false statements concerning events.


> Amazing



\section{Apuntes historia de la ciencia}

Esto es del artículo de Corina Yturbe %\cite{Yturbe1995}

Archimedes was probably the first mathematician candid enough to explain that
there is a difference between the way theorems are discovered and the way they
are proved.

ne of the philosophical theses in favor of the contraposition between the
internalist approach and the externalist approach is the so-called doctrine of
the two contexts, developed by positivism.

became the only aspect relevant to the history of science. The consideration of
non-rational, external factors, according to this position, not only fails to
increase our understanding of scientific development, but even obscures the
fundamental question of the rational validation of scientific arguments. p.75

The search for new explanatory programs is characterized above all by the
attempt to reconcile the internalist and externalist approaches. But, in view of
the fact that both approaches are committed to theses concerning the nature of
science that are not only incompatible, but are unsustainable, their union
without important changes in their philosophical presuppositions cannot result
in a viable program. p. 79

The general tendency in the historiography of the sciences is to consider the
social function of science, as weIl as so me aspects of the matrix from wh ich
the problematic is formed, as external elements, while the conceptual apparatus
and the problem field are treated as internal. But we should not think of
science as having two independent histories.

External factors are not only found in the context of discovery, they are
present also in the development of the concepts, problems, methods, problem
fields, etc. of scientific discourses: that is, external factors are present in
the context of validation itself. There are scientific discourses in which
ideological conceptions pass on to form part of the body of the science itself,
functioning as principles which define its field of study or guide its research;
thus, external factors can become internal. p. 85

This conception constitutes the dominant framework in which the philosophy of
science has been developed, and consists in drawing a radlcal distinction
between the context of discovery and the context of validation. p. 75

# Esto es del de Stillwell

This as certain risks, such as making the mathematics look simpler than it
really was in its time, but the risk of obscuring ideas by cumbersome,
unfamiliar notation is greater, in my opinion. Indeed, it is practically a
truism that mathematical ideas generally arise before there is notation or
language to express them clearly, and that ideas are implicit before they become
explicit. Thus the historian, who is presumably trying to be both clear and
explicit, often has no choice but to be anachronistic when tracing the origins
of ideas.

Perhaps the low cultural status of mathematics today, not to mention the
mathematical ignorance of politicians and philosophers, reflects the lack of an
Elements suitable for the modern world.

When mathematicians returned to problems of finding limits in the 17th century,
they found no use for the rigorous methods of the Greeks. The dubious
17th-century methods of infinitesimals were criticized by the Zeno of the time,
Bishop Berkeley, but little was done to meet his objections until much later,
since infinitesimals did not seem to lead to incorrect results. It was Dedekind,
Weierstrass, and others in the 19th century who eventually restored Greek
standards of rigor.

The story of rigor lost and rigor regained took an amazing turn when a
previously unknown manuscript of Archimedes, The Method, was dis- covered in
1906. In it he reveals that his deepest results were found using dubious
infinitary arguments, and only later proved rigorously. Because, as he says, “It
is of course easier to supply the proof when we have previously acquired some
knowledge of the questions by the method, than it is to find it without any
previous knowledge.” The importance of this statement goes beyond its revelation
that infinity can be used to discover results that are not initially accessible
to logic. Archimedes was probably the first  mathematician candid enough to
explainthat there is a difference between the way theorems are discovered and
the way they are proved.



# Esto es del de Hoyningen 15-07-2024

# Esto es del de Reichenbach. 19-07-2024

Thus a fictive definitive system of knowledge was made the basis of
epistemological inquiry, with the result that the schematized character of this
basis was soon forgotten, and the fictive construction was identified with the
actual system. It is one of the elementary laws of approximative procedure that
the consequences drawn from a schematized conception do not hold outside the
limits of the approximation; that in particular no consequences may be drawn
from features belonging to the nature of the schematization only and not to the
coordinated object.

for many a purpose the number tt maybe sufficiently approximated by the value
22/7; to infer from this, however, that tt is a rational number is by no means
permissible. Many of the inferences of traditional epistemology and of
positivism as well, I must confess, do not appear much better to me. It is
particularly the domain of the verifiability conception of meaning and of
questions connected with it, such as the problem of the existence of external
things, which has been overrun with paralogisms of this type.


# Esto es de Philipp Frank

The final theory has to be in fair agreement with observations and also has to
be sufficiently simple to be usable. If we consider this point, it is obvious
that such a theory cannot be "the truth." In modern science, a theory is
regarded as an instrument that serves toward some definite purpose. It has to be
helpful in predicting future observable facts on the basis of facts that have
been observed in the past and the present. It should also be helpful in the
construction of machines and devices that can save us time and labor. A
scientific theory is : in a sense, a tool that produces other tools according to
a practical scheme.

In Corder to obtain a basis from which one can pronounce sound judgment about
this situation, one should put the question: In what sense does science search
for the "truth" about the universe? This "truth" certainly does not consist of
"facts" but of general hypotheses or theories.

What we call "facts" in the strictly empirical sense, without any admixture of
theory, are, in the last analysis, sense impressions between which no connection
is given.

Hence, the question of what we have to regard as the "truth" about the universe
has to be formulated rather as follows: "What are the criteria under which we
accept a hypothesis or a theory?"

If we put this or a similar question, we shall see soon that these criteria will
contain, to a certain extent, the psychological and sociological characteristics
of the scientist, because they are relevant for the acceptance of any doctrine.

In other words, the validation of "theories" cannot be separated neatly from the
values which the scientist accepts. This is true in all fields of science, over
the whole spectrum ranging from geometry and mechanics to psychoanalysis.


sical conception, the aim of science is the discovery of an existing 'true'
world, while according to the positivist, however, it is the  construction of a
system of statements with the help of which we can  find our way in the world of
our experiences. Planck finds fault with  the latter conception: the passion and
readiness for sacrifice with  which men like Galileo have fought for their
convictions could not be  understood if the matter had been merely purposeful,
useful, constructions and not the discovery of truth.

The events around Galileo make it clear that the passionate conflicts  connected
with a physical theory have nothing to do with its suitability  to represent
natural processes but much more with their relationships  to the political and
social events of the time. Therefore there is no need  to amplify the positivist
conception of science by a metaphysical  concept of truth but only by a more
comprehensive study of the  connections that exist between the activity of the
invention of theories  and the other normal human activities.


. In this, I kept mainly to the formulations of Russian Marxists because the
connection between theory and  political practice is closest in their case. (p.
15) the law of causality and its limits Fuck me

arly years of its existence appears in Chapter III. Here we find the synthesis
of positivism  and the new logic explicitly presented. It is also  interesting
that since the rigorously logical formulation of the positivistic ideas, their
connection with  American pragmatism has become clearly revealed; in  this essay
this connection is distinctly emphasized.


In particular, Charles  W . Morris of Chicago recognized the connection  with
American pragmatism and publicized the idea  of cooperation between the two
groups. It was decided, for the purpose of this cooperation, to call a  special
congress, for which the name "Congress for  the Unity of Science" was coined by
Otto Neurath.  A preliminary conference in preparation for this congress took
place in 1934 in Prague. C. W . Morris said  there: "The doctrine of the Vienna
Circle is 'logical  positivism,' that of the American pragmatists is 'biological
positivism'." For the result of the cooperation  of the two, Morris suggested
the name "logical empiricism," which has been generally adopted in the  United
States.


Obviously, fitnessto support a desirable conduct of citizens or, briefly,to
support moral behavior, has served through the ages as a reason for acceptance
of a theory. When the "scientific criterions" did not uniquely determine a
theory, its fitness to support moral or political indoctrination became an
important factor for its acceptance. It is important to learn that the
interpretationof a scientifictheoryas a support of moral rules is not a rare
case but has played a role in all periods of history.

This book aims to treat all questions put in it without this opium  that is
often called 'philosophy', by misuse of an honourable name.  Nowhere will we
attempt to make unsettled questions appear as if they  were settled, by the use
of dazzling turns of phrase; moreover,  nowhere will we attempt to avoid
questions concerning the borders of  the area commonly called 'science', by
shifting to another plane. Much  more, our object is to treat all problems that
appear during the  operation of science with the same striving toward real
solutions to  which one is accustomed within the fields called often with benign
condescension the 'special sciences'.

. Nevertheless, many educators and political leaders were afraid that denial of
the exceptional status of the celestial bodies in physical science would make it
more difficult to teach the belief in the existence of spiritual beings as
distinct from material bodies; and since it was their general conviction that
the belief in spiritual beings is a powerful instrument to bring about a
desirable conduct among citizens, a physical theory that supported this belief
seemed to be highly desirable.

l these factors participate in the making of a scientific theory. We remember,
however, that according to the opinion of the majority of active scientists,
these extrascientific factors "should not" have any influence on the acceptance
of a scientific theory. But who has claimed and who can claim that they "should
not"?

. Technological changes have to produce social changes that manifest themselves
in changing human conduct. Everybody knows of the Industrial Revolution of the
nineteenth century and the accompanying changes in human life from a rural into
an urban patten^ Probably the rise of atomic power will produce analogous
changaf in man's way of life.


* Esto es algo que claramente tienen en mente los autores de "Not in our genes".


# Esto es del artículo de Salmon 'Statistical Explanation'

If inductive logic contains rules of inference which en ab le us to draw
conclusions from premises-much as in dedu ctive logi~then there is presumably
some number r which constitutes a lower bound for acceptance.




The problem of distinguishing between science and pseudoscience is part of the
larger task of determining which beliefs are epistemically warranted. This entry
clarifies the specific nature of pseudoscience in relation to other categories
of non-scientific doctrines and practices, including science denial(ism) and
resistance to facts. The major proposed criteria for defining and identifying
pseudoscience are discussed and some of their weaknesses are pointed out. There
is much more agreement on particular cases than on the general criteria that
such judgments should be based upon. This is an indication that there is still
much important philosophical work to be done on the relation between science and
pseudoscience.

We can have both theoretical and practical reasons for distinguishing between
real and false science (Mahner 2007, 516). From a theoretical point of view,
this distinction is an illuminating perspective that contributes to the
philosophy of science in much the same way that the study of fallacies
contributes to our knowledge of informal logic and rational argumentation. From
a practical point of view, the distinction is important for decision guidance in
both private and public life. Since science is our most reliable source of
knowledge in a wide range of areas, we need to distinguish scientific knowledge
from statements that are falsely claimed to be scientific.


here are many areas in which reliance on such statements can have disastrous
consequences:

Climate policy: The scientific consensus on ongoing anthropogenic climate change
leaves no room for reasonable doubt (Cook et al. 2016; Powell 2019). Science
denial has considerably delayed climate action, and it is still one of the major
factors that impede efficient measures to reduce climate change (Oreskes and
Conway 2010; Lewandowsky et al. 2019). Decision-makers and the public need to
know how to distinguish between competent climate science and science-mimicking
disinformation on the climate.

Healthcare: Medical science develops and evaluates treatments according to
evidence of their effectiveness and safety. Pseudoscience in healthcare gives
rise to ineffective and sometimes dangerous interventions and often lures people
away from science-based healthcare. Pseudoscience in preventive medicine makes
its victims refrain from vaccination and other efficient means to reduce the
risk of disease. Healthcare providers, insurers, government authorities and –
most importantly – patients and the general public need guidance on how to
distinguish between medical science and medical pseudoscience.

Environmental policies: In order to be on the safe side against potential
disasters it may be legitimate to take preventive measures when there is valid
but yet insufficient evidence of an environmental hazard. This must be
distinguished from taking measures against an alleged hazard for which there is
no valid evidence at all. Therefore, decision-makers in environmental policy
must be able to distinguish between scientific and pseudoscientific claims.

Expert testimony: It is essential for the rule of law that courts get the facts
right. The reliability of different types of evidence must be correctly
determined, and expert testimony must be based on the best available knowledge.
Sometimes it is in the interest of litigants to present non-scientific claims as
solid science. Therefore courts must be able to distinguish between science and
pseudoscience. Philosophers can contribute efficiently to the defence of science
against pseudoscience in such contexts (Pennock 2011; Ruse 2021).

Science education: The promoters of some pseudosciences (notably creationism)
try to introduce their teachings in school curricula. Teachers and school
authorities need to have clear criteria of inclusion that protect students
against unreliable and disproved teachings.

Journalism: When there is scientific uncertainty, or relevant disagreement in
the scientific community, this should be covered and explained in media reports
on the issues in question. Equally importantly, differences of opinion between
on the one hand legitimate scientific experts and on the other hand proponents
of scientifically unsubstantiated claims should be described as what they are.
Public understanding of topics such as climate change and vaccination has been
considerably hampered by organised campaigns that succeeded in making media
portray standpoints that have been thoroughly disproved in science as legitimate
scientific standpoints (Boykoff and Boykoff 2004; Boykoff 2008). The media need
tools and practices to distinguish between legitimate scientific controversies
and attempts to peddle pseudoscientific claims as science.

The common usage of the term “science” can be described as partly descriptive,
partly normative. When an activity is recognized as science this usually
involves an acknowledgement that it has a positive role in our strivings for
knowledge. On the other hand, the concept of science has been formed through a
historical process, and many contingencies influence what we call and do not
call science. Whether we call a claim, doctrine, or discipline “scientific”
depends both on its subject area and its epistemic qualities. The former part of
the delimitation is largely conventional, whereas the latter is highly
normative, and closely connected with fundamental epistemological and
metaphysical issues.

from here

@InCollection{sep-pseudo-science, author       =	{Hansson, Sven Ove}, title
=	{{Science and Pseudo-Science}}, booktitle    =	{The {Stanford} Encyclopedia
of Philosophy}, editor       =	{Edward N. Zalta and Uri Nodelman}, howpublished
=	{\url{https://plato.stanford.edu/archives/fall2025/entries/pseudo-science/}},
year         =	{2025}, edition      =	{{F}all 2025}, publisher    =
{Metaphysics Research Lab, Stanford University} }


This metaphysical picture quickly led to empiricist scruples, voiced by Berkeley
and Hume. If all knowledge must be traced to the senses, how can we have reason
to believe scientific theories, given that reality lies behind the appearances
(hidden by a veil of perception)? Indeed, if all content must be traced to the
senses, how can we even understand such theories? The new science seems to
postulate “hidden” causal powers without a legitimate epistemological or
semantic grounding. A central problem for empiricists becomes that of drawing a
line between objectionable metaphysics and legitimate science (portions of which
seem to be as removed from experience as metaphysics seems to be).


In the 1970s, a particularly strong form of scientific realism was advocated by
Putnam, Boyd, and others. When scientific realism is mentioned in the
literature, usually some version of this is intended. It is often characterized
in terms of these commitments:

Science aims to give a literally true account of the world. To accept a theory
is to believe it is (approximately) true. There is a determinate
mind-independent and language-independent world. Theories are literally true
(when they are) partly because their concepts “latch on to” or correspond to
real properties (natural kinds, and the like) that causally underpin successful
usage of the concepts. The progress of science asymptotically converges on a
true account.

Arguing that there is no fact of the matter about the geometry of physical
space. Poincaré proposed conventionalism: we decide conventionally that geometry
is Euclidean, forces are Newtonian, light travels in Euclidean straight lines,
and we see if experimental results will fit those conventions. Conventionalism
is not an “anything-goes” doctrine—not all stipulations will accommodate the
evidence—it is the claim that the physical meaning of measurements and evidence
is determined by conventionally adopted frameworks. Measurements of lines and
angles typically rely on the hypothesis that light travels shortest paths. But
this lacks physical meaning unless we decide whether shortest paths are
Euclidean or non-Euclidean. These conventions cannot be experimentally refuted
or confirmed since experiments only have physical meaning relative to them.
Which group of conventions we adopt depends on pragmatic factors: other things
being equal, we choose conventions that make physics simpler, more tractable,
more familiar, and so forth. Poincaré, for example, held that, because of its
simplicity, we would never give up Euclidean geometry. Scientific Realism and
Antirealism


The aim of science is: economy of thought (science is a useful instrument
without literal significance (Mach 1893)), the discovery of real relations
between hidden entities underlying the phenomena (Poincaré 1913), and the
discovery of a “natural classification” of the phenomena (a mathematical
organization of the phenomena that is the reflection of a hidden ontological
order (Duhem 1991)). These affinities, between 19th century global antirealism
and 20th century antirealism, mask fundamental differences. The former is driven
by methodological considerations concerning the proper way to do physics whereas
the latter is driven by traditional metaphysical or epistemological concerns
(about the meaningfulness and credibility of claims about goings-on behind the
veil of appearances).

”. For Duhem, epistemological holism holds only for physical theories for rather
special reasons; it does not extend to mathematics or logic and is not connected
with theses about meaning. Quine extends epistemological holism from physics to
all knowledge, including all knowledge traditionally regarded as a priori,
including allegedly analytic statements.] Quine, but not Duhem, believed that
our reluctance to revise mathematics and logic (because of their centrality to
our belief-systems) does not entail their a prioricity (irrevisability based on
evidence).

Moreover, if the analytic-synthetic distinction collapses, so too does the
positivist separation of metaphysics from science. For Quine, metaphysical
questions are just the most general and abstract questions we ask and are
decided on the grounds we use to decide whether electrons exist. All questions
are “internal” in the sense that they must be formulated in our home language
and answered with our standard procedures for gathering and weighing evidence.
In particular, questions about the reality of some putative objects are to be
answered in terms of whether they contribute to a useful organization of
experience and whether they withstand the test of experience.

Second, if T-terms were epistemologically or semantically problematic, that
would have to be due to the unobservable nature of their referents.


5. Scientific Realism

First, O-terms apply to apparently theoretical entities (for example, red
corpuscle) and T-terms apply to apparently observable entities (for example, the
moon is a satellite). Second, if T-terms were epistemologically or semantically
problematic, that would have to be due to the unobservable nature of their
referents. But in the continuous gradation between seeing with the unaided eye,
with binoculars, with an optical microscope, with an electron microscope, and so
on, there is no sharp cut-off between being observable and being unobservable
where we could non-arbitrarily say: beyond this we cannot trust the evidence of
our senses or apply terms with confidence.

In the 1970s, a particularly strong form of scientific realism (SR) was
advocated by Putnam, Boyd, and others (Boyd 1973, 1983; Putnam 1962, 1975a,
1975b). When scientific realism is mentioned in the literature, usually some
version of SR is intended. SR is often characterized in terms of two commitments
(van Fraassen 1980):

SR1     Science aims to give a literally true account of the world.

SR2     To accept a theory is to believe it is (approximately) true.

However, scientific realists’ arguments and their interpretation of SR1 and SR2
often presuppose further commitments:

SR3     There is a determinate mind-independent and language-independent world.

SR4     Theories are literally true (when they are) partly because their
concepts “latch on to” or correspond to real properties (natural kinds, and the
like) that causally underpin successful usage of the concepts.

SR5     The progress of science asymptotically converges on a true account.



The positivist may respond that they cannot be directly sensed, and are thus
unobservable, but why should being directly sensed be the criterion for
epistemological or semantic confidence? Fourth, observation is theory-infected:
what we can both observe and employ as evidence is a function of the language,
concepts, and theories we possess. A primitive Amazonian may observe a tennis
ball (he notices it), but without the relevant concepts he cannot use it as
evidence for any claims about tennis. Such arguments undermine a central
distinction of the positivist program.

First, a few clarifications of IBE are in order. If IBE is to be non-trivial,
the best explanation must not entail that what is best must antecedently be what
is most likely, since of course we should infer the truth of the most likely
explanation. Rather the best explanation must be characterized in terms of
properties like “loveliest” or “most explainey” (Lipton 2004). Traditional
examples of such properties are: it has wide scope and precision; it appeals to
plausible mechanisms; it is simple, smooth, elegant, and non-ad hoc; and it
underwrites contrasts (why this rather than that). Then IBE says we should
accept the theory that optimizes such explanatory virtues when explaining the
phenomena. The caveat “if there is one” blocks inferences to the best of a bad
lot: the best explanation may not reach a minimally acceptable threshold.
Finally, like any inferential principle that amplifies our knowledge,
conclusions inferred by IBE are fallible: while they are more likely to be true,
they could be false. Second, the “justification” for IBE is two-fold. (1) It is
needed for science. Simple enumerative induction (which entitles us to move
probabilistically from “All observed As are Bs” to “All As are Bs” cannot handle
inferences from observed phenomena to their “hidden” causes. For example, we
cannot inductively infer “Galaxy X is receding” from “Light from Galaxy X is
red-shifted”, but we can infer by IBE that Galaxy X is receding because that is
the best explanation of why its light is red-shifted. More strongly, Harman
(1965) argues that IBE is needed to warrant straight enumerative induction: we
are entitled to make the induction from “All observed As are Bs” to “All As are
Bs” only if “All As are Bs” provides the best explanation of our total evidence.
(2) Scientific uses of IBE are grounded in, and are just sophisticated
applications of, a principle we use in everyday inferential practice. If I see
nibbled cheese and little black deposits in my kitchen and hear scratching
noises in the walls, I reasonably infer that I have mice, because that best
explains my evidence. IBE thus needs no more justification than does modus
ponens—each is part of the very practices that constitute what rational
inference is.


\section{Constructive Empiricism}

Van Fraassen (1980) proposed constructive empiricism (CE), arguing that we can
preserve the epistemological spirit of positivism without subscribing to its
letter. Van Fraassen’s is an antirealism concerning unobservable entities.
Recognizing the difficulties of basing antirealism on a “broken-backed”
linguistic distinction between O-terms and T-terms, he allows our judgments
about unobservables to be literally construed but, he argues, our evidence can
never entitle us to our beliefs about unobservables. CE is consistent with SR3
and SR4 (though it does not commit to them, it has no quarrels with realist
objectivity or semantics) but replaces SR1, SR2, and SR5 respectively with:

CE1     Science aims to provide empirically adequate theories of the phenomena.

CE2     To accept a theory is to believe it is empirically adequate, but acceptance has further non-epistemic/pragmatic features.

CE5     The progress of science produces increasing empirical adequacy.

A theory T is empirically adequate if and only if what T says about all actual
observable things and events is true (that is, T saves all the phenomena, or T
has a model that all actual phenomena fit in). Empirical adequacy is logically
weaker than truth: T’s truth entails its empirical adequacy but not conversely.
But it is still quite strong: an empirically adequate theory must correctly
represent all the phenomena, both observed and unobserved. CE2 distinguishes
epistemic and pragmatic aspects of acceptance. Epistemic acceptance is belief;
beliefs are either true or false. Pragmatic acceptance involves non-epistemic
commitments to use the theory in certain ways (basing research, experiments,
and explanations on it, for example); commitments are neither true nor false;
they are either vindicated or not. CE5 acknowledges that there is instrumental
progress without trying to explain it. CE concedes a realist semantics
(“electron”-talk is not highly derived talk about observables) but preserves
the spirit of positivism by recommending agnosticism about a theory’s literal
claims about unobservables.

\subsection{a. The Semantic View of Theories and Empirical Adequacy}

On the positivist view, a theory T is a syntactic object: T is the set of
theorems in a language generated from a set of axioms (the laws of T) and
derivation rules. The empirical content (the entire literal content) of T is
T/O, the theorems expressible in the observational vocabulary. A theory T is
empirically (observationally) adequate if T/O is the class of all true
observational sentences.  Two theories, T and T’, are empirically
(observationally) equivalent if T/O = T’/O. Since such theory pairs have the
same literal content and differ only in their non-literal, theoretical content,
they are merely inter-definable variants of a common observational basis: they
say the same thing but express it differently. There is no fact of the matter
whether T or T’ is true (both are or neither are), and whether we work with T
or T’ is purely a pragmatic matter concerning which is simpler, more
convenient, and so forth. For SR and CE there is a fact of the matter: at most
one of T, T’ can be true. For SR there may be reasons to believe one of T, T’.
For CE there can be no epistemic reason to believe one over the other, though
there may be pragmatic reasons to accept (commit to using) one over the other.
Van Fraassen needs a different account of theories if he is to agree with
realists about literal content and there being a fact of the matter about
empirically equivalent theories.

For him, a theory T is a semantic object, the class of models, A = <D, R1, R2,
…, Rn>, that satisfy its laws (where D is a set of objects and Ri are
properties and relations defined on them). For example, D might contain
billiards and molecules; the property is elastic in A might be instantiated by
both billiards and molecules, is a molecule by some members of D, and is a
billiard ball by others. Now let A’ = <D’, R’1, R’2, …, R’m> (where m < n, D’
is a proper subset of D, and R’i = Ri/D’ (Ri restricted to D’)). Intuitively A’
is obtained from A by removing all unobservables, so D’ would contain billiard
balls but not molecules, is elastic would now be restricted to billiard balls,
is a molecule would not be instantiated, and so forth. Then A’ is an empirical
substructure of A, the result of restricting the original domain to observables
and its properties and relations accordingly. T is empirically adequate if and
only if T has an empirical substructure that all observables fit in. Two
theories, T and T’, are empirically equivalent if all the observables in a
model of T are isomorphic to the observables in a model of T’. Such theory
pairs agree in what they say about observables but may disagree in what they
say about unobservables. Thus CE can agree with SR that at most one of T, T’
can be true and to be a realist about that theory is to believe it is true
(SR2). Yet CE can preserve the spirit of positivism by holding that we can
never have reason to believe a theory; at most we have reason to believe it is
empirically adequate. Friedman (1982) questions whether van Fraassen achieves
this.


Michael Liston Email: mnliston@uwm.edu University of Wisconsin-Milwaukee U. S.
A.
